\section{第二阶段深讲：调度器工程化，从算法到多核策略}

调度算法课常用 FCFS、RR、优先级、MLFQ、CFS、EDF 做比较，但内核调度器不是算法博物馆。真实调度器要在中断、锁、cache、NUMA、cgroup、实时任务、能耗和安全边界之间做决策。本节关注算法落地后的工程问题：运行队列怎么组织，唤醒如何抢占，负载如何迁移，配额如何限流，证据如何解释。

\subsection{调度器只从 runnable 集合里选}

很多性能分析误把“线程没运行”归因于调度器。调度器只负责选择 runnable 任务；如果任务 sleeping，它等待的是锁、I/O、page fault、timer、信号或条件变量。正确诊断必须先问任务在哪个集合中。

\begin{longtable}{p{0.2\linewidth}p{0.36\linewidth}p{0.32\linewidth}}
\toprule
任务位置 & 含义 & 调度器能否直接运行 \\
\midrule
CPU current & 正在运行或在内核入口中 & 已经在运行 \\
runqueue & runnable，等待 CPU & 可以被选择 \\
wait queue & 条件不满足，睡眠中 & 不能，必须被唤醒 \\
completion/request & 等设备或异步结果 & 不能，completion 路径唤醒 \\
zombie list & 已退出，等父进程 wait & 不能 \\
\bottomrule
\end{longtable}

这个区分决定了工具选择。runnable wait 高，查 runqueue 和 CPU 配额；sleep wait 高，查锁、I/O 或等待对象；zombie 多，查父进程是否 wait。没有这个分类，调度分析会变成猜测。

\subsection{每 CPU runqueue 为什么必要}

单一全局 runqueue 实现简单，但多核上所有 CPU 都要竞争一把锁，且每次调度都会移动同一组 cache line。每 CPU runqueue 把局部调度变成局部锁，降低争用；代价是负载不均时需要跨 CPU 迁移。

\begin{lstlisting}
per-cpu scheduling:
  local timer tick updates current CPU's rq
  local wakeup usually enqueues on target CPU rq
  idle CPU tries to pull tasks from busy CPU
  periodic balance compares load across domains
\end{lstlisting}

这就产生了调度域：SMT sibling、同一 LLC、同一 NUMA node、跨 node。近距离迁移成本低，远距离迁移成本高。调度器不能只看任务数量，还要看任务负载、cache 热度、内存位置和 CPU 能力。

\subsection{唤醒路径比定时器路径更影响交互延迟}

交互任务多数时间在睡眠，事件到来后被唤醒。如果唤醒后还要等很久才运行，用户会感到卡顿。因此调度器不仅要在时间片耗尽时选择任务，还要在 wakeup 时判断是否应抢占当前任务。

\begin{lstlisting}
wakeup(task):
  task.state = RUNNABLE
  place task on target runqueue
  if task should preempt current:
    set reschedule flag
    if remote CPU:
      send reschedule IPI
\end{lstlisting}

这里有两个难点。第一，唤醒任务放到哪个 CPU：放回上次 CPU 有 cache 局部性，放到空闲 CPU 有低延迟，放到数据所在 NUMA node 有内存局部性。第二，是否抢占当前任务：过度抢占会增加切换成本，抢占不足会增加交互延迟。

\subsection{CFS 的 \code{min\_vruntime} 是队列坐标系}

只讲 \code{vruntime} 还不够。每个 CFS runqueue 还维护 \code{min\_vruntime}，表示队列已经推进到的公平时间基准。新任务和唤醒任务不能简单设为 0，否则它们会在树最左边占便宜；也不能放得太靠后，否则唤醒延迟差。

\begin{lstlisting}
enqueue_fair(task, rq):
  if task is new:
    task.vruntime = rq.min_vruntime
  if task wakes from sleep:
    task.vruntime = max(task.vruntime, rq.min_vruntime - wakeup_granularity)
  insert task into rb_tree ordered by vruntime
\end{lstlisting}

\code{min\_vruntime} 让每个 runqueue 有自己的坐标系。迁移任务时，需要把任务的 \code{vruntime} 从源 runqueue 坐标转换到目标 runqueue 坐标，否则任务可能因为跨 CPU 迁移得到不公平优势或惩罚。

\subsection{EEVDF 思路：从公平份额到最早合格截止}

现代公平调度进一步引入虚拟 deadline 思路。直观地说，任务不仅有“已经用了多少 CPU”的账本，还可以有“按权重和 slice 计算出的虚拟截止时间”。调度器倾向选择已经 eligible 且虚拟 deadline 最早的实体，以改善延迟和公平之间的折中。

\begin{lstlisting}
entity fields:
  vruntime
  weight
  slice
  virtual_deadline = vruntime + slice / weight_scaled

pick:
  among eligible entities
  choose earliest virtual_deadline
\end{lstlisting}

这里不需要把工业实现细节全部搬进原理卷，但要说明动机：单纯最小 \code{vruntime} 关注长期公平，交互延迟还要通过 slice、wakeup 和 eligibility 调节。调度器演进的方向，是把“公平”和“延迟目标”放进同一个可维护账本。

\subsection{实时调度和普通调度的边界}

实时任务不能和普通 CFS 任务只靠 nice 值竞争。实时调度类通常优先于普通调度类，FIFO/RR 实时任务能长时间压制普通任务。deadline 调度还要处理 runtime、period、deadline 和准入控制。

\begin{longtable}{p{0.2\linewidth}p{0.34\linewidth}p{0.34\linewidth}}
\toprule
调度类 & 目标 & 风险 \\
\midrule
stop/deadline & CPU 停止、deadline 保证 & 最高优先级路径错误会冻结系统 \\
real-time FIFO/RR & 固定优先级实时响应 & 饥饿普通任务，需要限额 \\
fair & 权重公平与通用吞吐 & 不能保证硬实时 \\
idle & 没有其他任务时运行 & 不应影响正常调度 \\
\bottomrule
\end{longtable}

实时任务需要限额和监控，否则一个死循环实时线程可以让系统失去响应。工业内核通常提供 RT throttling，防止实时任务占满全部 CPU 时间。

\subsection{cgroup 配额造成的周期性延迟}

容器或服务组常被配置 CPU quota。假设周期 100ms，配额 20ms。组内任务在本周期累计运行 20ms 后会被 throttle，即使机器还有空闲 CPU，也要等下个周期补充预算。

\begin{lstlisting}
cgroup runtime:
  on task execution:
    group.runtime_remaining -= delta
  if group.runtime_remaining <= 0:
    throttle all runnable tasks in group
  at period timer:
    refill group.runtime_remaining
    unthrottle group
\end{lstlisting}

这种机制会形成很典型的尾延迟：请求前 20ms 很快，预算耗尽后突然等待到周期刷新。排查容器内延迟时，必须看 throttled time、throttled periods 和 quota 配置，而不是只看宿主 CPU 利用率。

\subsection{优先级反转是调度器和锁的交叉问题}

低优先级任务持锁，高优先级任务等锁，中优先级任务占 CPU，会导致高优先级任务间接被中优先级任务阻塞。调度器单独看 runnable 集合时看不出问题，因为真正持有资源的是低优先级任务。

\begin{lstlisting}
priority inheritance:
  H waits on lock owned by L
  boost L effective priority to at least H
  if L waits on another owner X:
    propagate boost to X
  when lock released:
    recompute effective priority
\end{lstlisting}

实现上要维护 owner、waiters、effective priority 和继承链。释放锁时不能简单恢复原优先级，因为任务可能还持有另一个被高优先级任务等待的锁。这个例子说明同步原语不能和调度器完全隔离。

\subsection{调度器证据：如何解释一次延迟}

一次请求慢，要按时间分段：

\begin{lstlisting}
request latency =
  runnable wait
  + on-cpu execution
  + sleeping on I/O or lock
  + page fault/reclaim
  + scheduler and interrupt overhead
\end{lstlisting}

常用证据包括 \code{perf sched}、\code{/proc/<pid>/sched}、tracepoint、off-CPU profile、\code{/proc/schedstat}、cgroup CPU 统计、\code{pidstat -w}。不同证据回答不同问题。

\begin{longtable}{p{0.24\linewidth}p{0.46\linewidth}p{0.18\linewidth}}
\toprule
证据 & 回答的问题 & 层次 \\
\midrule
\code{/proc/<pid>/sched} & 任务运行、等待、迁移等统计 & 单任务 \\
\code{perf sched} & 调度事件时间线 & 系统级 \\
cgroup cpu.stat & 是否被 quota throttle & 服务组 \\
tracepoint sched\_* & wakeup、switch、migrate 事件 & 内核事件 \\
off-CPU profile & 睡眠等待在哪里发生 & 锁/I/O \\
\bottomrule
\end{longtable}

调度器优化要先定位是哪一段时间异常。若大部分时间在 off-CPU 等锁，改 CFS 参数不会解决；若大部分时间 runnable wait，才应看 CPU 配额、优先级、runqueue 和负载均衡。

\subsection{调度器工程验收}

\begin{itemize}
\tightlist
\item 能解释 task 为什么不在 CPU：runnable、sleeping、zombie、throttled 要分清。
\item 能说明每 CPU runqueue 与负载均衡的取舍。
\item 能推导唤醒抢占对交互延迟和上下文切换成本的影响。
\item 能解释 \code{vruntime}、\code{min\_vruntime} 和迁移坐标转换的关系。
\item 能说明 cgroup quota 为什么会制造周期性尾延迟。
\item 能用调度事件证据区分 on-CPU 慢和 off-CPU 等待。
\end{itemize}
